\documentclass[conference]{IEEEtran}
\usepackage[utf8]{inputenc}
\usepackage[T1]{fontenc}
\usepackage{cite}
\usepackage{amsmath,amssymb}
\usepackage{graphicx}
\usepackage{booktabs}
\usepackage{url}

\title{Effectiveness of a Generate--Verify--Repair Pipeline with Batfish for LLM‑Generated Vendor CLI Configurations}

\author{
\IEEEauthorblockN{Autonomous AI Researcher}
\IEEEauthorblockA{CNSM 2026 Agentic AI Researcher Track}
}

\begin{document}

\maketitle

\begin{abstract}
We report a preregistered paired experiment (cnsms2026-001-gvr-batfish-prereg-v1) that evaluated whether a generate--verify--repair (GVR) pipeline---shared initial LLM candidate (gpt-5-mini) + a deterministic Batfish‑style validator + at most one automated LLM repair---reduces deterministic‑verifier detected semantic/safety failures on intent\ensuremath{\rightarrow}vendor‑CLI tasks. Planned N=300 pairs; 266 complete pairs were analysed after preregistered exclusions. Baseline pass-rate = 260/266 (97.744\%); guarded (final GVR) pass-rate = 265/266 (99.624\%). Paired difference = 0.01879699248; exact McNemar two‑sided p = 0.0625. Paired bootstrap (B=10,000, seed=1729) 95\% percentile CI = [0.0037593984962406013, 0.03759398496240601]. Repair was invoked for 6 tasks; median repair iterations per task = 0. Per‑episode latency was not recorded in per‑task artifacts and thus could not be evaluated. Under shared\_initial\_candidate semantics the observed absolute improvement is small and did not meet the preregistered \ensuremath{\alpha}=0.05 decision rule. We provide preregistered design details, deterministic validator semantics, exact paired counts, representative artifact‑grounded repair examples, contamination findings (no exact public duplicates flagged), and a focused discussion of limitations and implications.
\end{abstract}

\section{Introduction}
Generative large language models (LLMs) are actively explored for NetOps tasks such as intent\ensuremath{\rightarrow}CLI translation, zero‑touch configuration synthesis, and automated maintenance workflows [1], and surveys emphasize both opportunity and risk when LLMs are applied to operational network configuration [2]. Stochastic LLM outputs can produce syntactic or semantic violations---missing required commands, unsafe command orderings, or constraint breaches---that create operational risk if applied directly. To mitigate that risk, pipelines that interpose deterministic verifiers (e.g., Batfish or header‑space analyses) and bounded repair loops---collectively generate--verify--repair (GVR)---have been proposed to provide deterministic acceptance criteria and automated correction [3,4]. Prior work presents components and prototypes, but preregistered, paired evaluations that archive verifier traces and quantify verifier‑triggered repair benefit on reproducible NetOps task banks remain limited.

This study tested the preregistered question (cnsms2026-001-gvr-batfish-prereg-v1): Does a GVR pipeline that uses a single sampled initial LLM candidate per task (shared\_initial\_candidate semantics), a deterministic Batfish‑style validator, and at most one automated LLM repair call reduce deterministic‑verifier detected semantic/safety failures relative to the same initial candidate without repair? Our principal contribution is a preregistered, artifact‑archived paired evaluation executed under the CNSM Agentic AI Researcher constraints that reports exact paired counts, inferential outcomes, execution accounting, representative artifact examples, and an evidence‑grounded interpretation.

\section{Related Work}
LLM applications in NetOps span intent‑to‑configuration translation, automated configuration generation, and AIOps toolchains; several recent benchmark and survey efforts document these applications and their limitations [1,2]. Task‑level benchmarks such as NetConfEval examine whether LLMs can synthesize vendor CLI fragments from intent descriptions and show that vendor‑specific syntax, sequencing requirements, and implicit ordering constraints remain failure points for unconstrained generation [1]. Broad surveys of LLMs for network operations synthesize common failure modes---hallucination, constraint violations, and brittle prompt dependence---and recommend grounding, verifier integration, and tooling to reduce operational risk [2].

Generate--verify--repair (GVR) architectures have been proposed in adjacent domains (e.g., code generation and regulated document production) to combine a stochastic generator with a deterministic oracle and bounded repair steps; these works articulate architectures and convergence desiderata for achieving reproducible, verifiable outputs from stochastic agents [3]. In NetOps, deterministic network verifiers (Batfish and related header‑space analyses) provide concrete, automatable checks for reachability, ordering, and policy constraints; Batfish's engineering history highlights tradeoffs in rule expressiveness and scalability that directly affect its suitability as a validator in automated loops [4]. Empirical and analytic studies of what LLMs require to synthesize correct router or device configurations emphasize that correct sequencing, vendor dialect handling, and explicit constraint encoding are necessary for high first‑pass correctness and motivate verifier‑driven repair when those elements are missing in a generated candidate [5].

Taken together, the literature establishes three relevant points for our design: (1) NetOps generation tasks combine syntactic, ordering, and topology‑dependent semantic constraints that are poorly captured by surface‑level language modeling alone [1,5]; (2) deterministic verifiers can provide a domain‑appropriate acceptance criterion but have finite rule coverage and impose engineering tradeoffs that shape experimental outcomes [4]; and (3) GVR‑style pipelines are a pragmatic mitigation pattern that can convert nondeterministic LLM outputs into verifier‑acceptable artifacts if the repair stage can correct verifier‑reported defects within bounded calls [3]. Our study differs from prior work by executing a preregistered paired test using shared\_initial\_candidate semantics, by archiving raw model and validator traces for auditability, and by reporting exact paired contingency counts together with the preregistered exact McNemar test and a paired bootstrap CI to help interpret small‑discordant outcomes.

\section{Methodology}
Overview and estimand.  We preregistered the study as cnsms2026-001-gvr-batfish-prereg-v1 and analysed the paired estimand labeled paired\_success\_rate\_difference\_guarded\_minus\_baseline under shared\_initial\_candidate semantics. Under those semantics the same single sampled initial model candidate (per task) is used to compute both outcomes: baseline = validator pass/fail on the initial candidate; guarded = final validator pass/fail after the allowed validator‑triggered repair (at most one LLM repair call). This design isolates the marginal effect of the permitted validator‑triggered repair for a fixed initial candidate rather than conflating effects of different initial samples or different first‑pass prompting strategies.

Task generation, inclusion/exclusion, and determinism.  The preregistered task bank deterministically produced 300 intent\ensuremath{\rightarrow}CLI pairs composed of public NetConfEval‑style examples and synthetic tasks generated by a seeded deterministic generator. Tasks are stratified into two vendor‑dialect clusters (150 Cisco‑like, 150 Juniper‑like). Generator IDs, seeds, and per‑task payloads (initial\_state, expected\_final\_state, required\_sequence, allowed\_touched\_fields, and workflow\_policies) are archived in the task manifest. Inclusion/exclusion rules and the retry policy followed preregistration: transient provider/API failures were retried up to two times; pairs were excluded from the primary confirmatory analysis only when both arms failed due to infrastructure/provider errors consistent with the preregistered missingness\_plan. All task selection indices were fixed before model outcomes were observed and are archived to ensure deterministic replay of the task bank.

Model, orchestration and recorded runtime parameters.  The declared model used in the run was gpt-5-mini (execution/model\_configuration.json). Archived execution metadata records max\_output\_tokens = 1024, maximum\_attempts\_per\_call = 1, provider = openai\_responses, and temperature = null (unset). Provider accounting in the deterministic reconciliation artifact records hosted\_model\_call\_event\_count = 306, completed\_hosted\_model\_call\_event\_count = 272, and failed\_hosted\_model\_call\_event\_count = 34; stage accounting shows shared\_initial\_generation = 300 and repair = 6. Because the provider configuration recorded temperature as unset and no centralized deterministic sampling seed is available for hosted calls, nondeterministic sampling of provider outputs cannot be excluded; we disclose this operational nondeterminism and treat it as a limitation in the Results and Disclosure.

Deterministic validator semantics and scoring rules.  The binary oracle was deterministic\_netops\_validator\_v5 and it applied a fixed battery of checks recorded in per‑episode validator traces. For each candidate the validator (1) parsed and normalized the CLI commands (producing parsed\_command\_count and normalized\_configuration), (2) checked the required\_sequence field against parsed commands (required\_command\_count), (3) enforced constrained\_touch policies (allowed\_fields and allowed\_touched\_fields), and (4) executed reachability/constraint checks on the synthetic topology model to detect sequencing or dependency violations. The validator output includes observed\_touched\_field\_count, violation\_codes, and final valid boolean; the primary endpoint was the validator's binary pass/fail under the scoring rule full\_intent\_constraint\_validation. Every episode archives both validation\_before and validation\_after traces so individual failure modes are inspectable and replayable.

Repair loop mechanics (bounded and archived).  The preregistered repair stage is strictly bounded to at most one automated LLM repair call per task when the initial candidate fails the validator. The repair prompt construction is deterministic given the validator failure codes and the task payload: it includes the task constraints, the normalized initial candidate, and explicit violation codes, and it requests a corrected configuration that satisfies the validator battery. Repair provider traces and repair\_prompt\_sha256 values are archived when invoked; stage accounting reports six repair calls in the run. Repair outputs are submitted to the same deterministic validator and the guarded final outcome is the validator result after that single repair attempt. Limiting repairs to one invocation preserves a clear estimand and enforces a bounded repair budget as preregistered.

Analysis procedures and exact inferential specification.  The prespecified primary inferential test is the exact McNemar two‑sided test on paired binary outcomes at \ensuremath{\alpha}=0.05. To quantify effect magnitude and uncertainty we also preregistered a paired bootstrap percentile confidence interval with B = 10,000 resamples and seed = 1729; those exact bootstrap parameters are recorded in analysis/analysis\_log.jsonl. Missingness handling followed the preregistered complete\_pair\_primary rule: pairs where both arms failed due to provider infrastructure issues were excluded from the primary test (documented in analysis/missingness\_summary.csv), and conservative sensitivity analyses treating excluded pairs as failures were preserved as archived artifacts.

Accounting, reconciliation and reproducibility artifacts.  Execution accounting recorded planned\_episode\_count = 600 and completed\_episode\_count = 532, with failed\_episode\_count = 68; after applying preregistered exclusions the analysed complete\_pair\_count = 266. Deterministic reconciliation artifacts verify internal consistency of paired counts and provider calls (provider and stage counts, cache behavior, and cross‑condition accounting) and assert that deterministic consistency checks passed. Key artifact categories archived for reproducibility include raw model responses, provider call traces, per‑episode scoring JSONs (validator traces), the task manifest, transformation manifests, model\_configuration.json, execution/raw\_results.jsonl, and the analysis result artifacts (condition\_summary.csv, paired\_contingency\_table.csv, missingness\_summary.csv, contamination\_summary.csv, and the paired\_difference\_figure.svg). File paths and manifests are provided with the execution bundle to enable exact reanalysis and targeted replay of repaired episodes.

Operational measurements attempted and absent.  The preregistration included secondary estimands for median\_repair\_iterations\_per\_task and median end‑to‑end latency; repair iterations were measurable and recorded (repair calls = 6; median repair iterations per task = 0). Per‑episode end‑to‑end latency (generation\ensuremath{\rightarrow}validation\ensuremath{\rightarrow}repair\ensuremath{\rightarrow}final validation) was not recorded in the archived per‑task artifacts and therefore the preregistered latency estimand could not be evaluated; this absence is noted as a preregistered deviation in the Limitations and drives a concrete reproducibility recommendation to include timing instrumentation in future runs.

Threats to validity embedded in the design.  The methodology preserves explicit distinctions required by the preregistration: internal validity is governed by shared\_initial\_candidate semantics (the estimand measures repair benefit for a fixed sample rather than differences from alternative first‑pass strategies), construct validity depends on validator rule coverage (the deterministic\_netops\_validator\_v5 battery defines which semantic violations are detectable), and external validity is limited by the two vendor clusters and synthetic topology scale. The preregistration's missingness rules and sensitivity analyses were applied exactly and all exclusions, retries, and provider failures are archived in the execution manifest to support transparent effect‑size interpretation.

\section{Results}
Execution accounting and sample flow: The preregistered plan targeted 300 independent intent\ensuremath{\rightarrow}CLI tasks (600 episodes across two conditions). The run recorded planned\_episode\_count = 600 and completed\_episode\_count = 532; provider/infrastructure transient failures produced failed\_episode\_count = 68. Applying the preregistered 'complete\_pair\_primary' exclusion rules (exclude both‑arm provider/infrastructure failures) yielded complete paired units analysed = 266. The analysis pipeline's provider accounting indicates hosted\_model\_call\_event\_count = 306 total model‑call events across the run, with completed\_hosted\_model\_call\_event\_count = 272 and failed\_hosted\_model\_call\_event\_count = 34; stage accounting records shared\_initial\_generation = 300 and repair = 6.

Primary paired outcomes and contingency counts: For the 266 complete pairs the validator outcomes produced the following contingency counts: n11 = 260 (both baseline and guarded pass), n10 = 5 (guarded pass only), n01 = 0 (baseline pass only), and n00 = 1 (both fail). These counts yield observed per‑condition proportions baseline = 260/266 = 0.9774436090 and guarded = 265/266 = 0.9962406015; the paired (guarded - baseline) difference = 0.01879699248. The exact McNemar two‑sided test applied to the discordant pair counts (n10 = 5, n01 = 0) produced p = 0.0625. Because the exact two‑sided McNemar p for observed discordants is 0.0625 (1/16), it narrowly misses the preregistered \ensuremath{\alpha} = 0.05 rejection threshold.

Interpreting the small‑discordant pattern and the two inferential summaries: The discordant total (n10 + n01 = 5) is small. The exact two‑sided McNemar test calculates two‑sided probability mass under a Binomial(total=5, p=0.5) for outcomes as or more extreme than the observed (k >= 5 or k <= 0), yielding p = 2 * (0.5\textasciicircum{}5) = 0.0625. In contrast, the paired bootstrap percentile CI (B = 10,000 resamples; seed = 1729) for the paired difference is [0.0037593985, 0.0375939850] and excludes zero. Both summaries are valid descriptions of uncertainty; the discrete exact test is conservative when discordant counts are few, while the bootstrap CI summarizes empirical resampling variability. We therefore report both: the preregistered hypothesis test (McNemar) did not reject at \ensuremath{\alpha} = 0.05, while the paired bootstrap CI suggests a small positive effect whose lower bound is barely above zero for the observed sample.

Repair invocation, yield, and derived operational quantities: Repair was invoked in 6 episodes (stage\_counts: repair = 6). The contingency pattern (n10 = 5) indicates that at most five of those repair invocations contributed to additional guarded successes among the 266 analysed pairs; the run-level accounting (6 repair calls vs 5 discordant guarded wins) is consistent with a small number of successful repairs and at most one repair that did not produce a guarded pass among the complete pairs. Median repair iterations per task across complete pairs is 0 (most tasks did not require repair). A useful derived metric is the number needed to treat (NNT) interpretation: 1 / absolute\_risk\_reduction = 1 / 0.01879699248 \ensuremath{\approx} 53.2, i.e., roughly 53 guarded tasks (under this run's conditions and shared\_initial\_candidate semantics) were associated with one additional verifier‑approved configuration relative to baseline. This NNT is a descriptive operational scalar and should be interpreted in context: with a very high baseline pass rate, the marginal correction yield per repair opportunity is small.

Representative repair case diagnostics: Representative artifacts illustrate typical failure and repair behavior. Task-000008 (difficulty: "extreme", required\_command\_count = 9) provides a concrete example: the shared initial candidate contained a truncated final token line ("interface uplink2\textbackslash{}n") and failed the validator's parsed\_command\_count vs required\_command\_count check; the automated repair call produced a corrected candidate that completed the previously truncated command and included the required edge1 sequence. The guarded final validator trace for task-000008 shows validation\_after.valid = true and violation\_count = 0. These per‑episode validator traces record parsed\_command\_count, required\_command\_count, observed\_touched\_field\_count, normalized\_configuration, and violation codes, enabling targeted audit of repair logic and verifier coverage for each repaired episode.

Sensitivity to missingness and preregistered conservative handling: Thirty‑four both‑arm episodes were excluded from the primary confirmatory analysis (complete\_pairs = 266) per preregistered rules. A conservative sensitivity analysis that treats those 34 excluded both‑arm episodes as failures (complete N = 300) yields baseline = 260/300 = 0.866667 and guarded = 265/300 = 0.883333; the discordant counts among the complete pairs remain n10 = 5, n01 = 0, producing the same exact McNemar p = 0.0625 for the paired test on the complete‑pair subset. This shows the primary inferential conclusion (nonrejection under exact McNemar) is robust to preregistered conservative missingness treatment, while the absolute per‑condition proportions change due to inclusion of excluded episodes as failures.

Quantitative artifact summaries and reproducibility metadata used in Results: All numerical summaries reported above derive directly from archived analysis artifacts: the paired contingency counts, condition summaries, missingness summary, and provider‑call accounting produced the reported totals (complete\_pairs = 266; baseline\_success\_count = 260; guarded\_success\_count = 265; hosted\_model\_call\_event\_count = 306; repair calls = 6; both\_missing\_pairs = 34). The paired bootstrap used B = 10,000 resamples with seed = 1729 (bootstrap parameters archived) to produce the reported percentile CI. Per‑episode validator traces and scoring JSONs include normalized\_configuration, parsed\_command\_count, required\_command\_count, and violation codes for each episode and can be used to reproduce failure‑mode breakdowns or to recalibrate validator rule coverage.

Limitations specific to the observed Results: (1) High baseline pass rate produced few discordant pairs: rare failures reduce the statistical power to detect small absolute improvements under the preregistered exact McNemar test. (2) Repair calls were infrequent (6/300) in this corpus; larger absolute improvements from GVR are more plausible on corpora with lower first‑pass accuracy. (3) The NNT and repair yield reported are conditional on shared\_initial\_candidate semantics (the same initial sample across arms) and the particular model sampling/settings used; different sampling regimes or constrained first‑pass prompts could change the yield. (4) Per‑episode latency was not archived and thus the preregistered latency estimand could not be evaluated; timing instrumentation is recommended for future runs. (5) Validator coverage is finite: violations outside deterministic\_netops\_validator\_v5's rule battery are unmeasured here.

\section{Discussion}
Statistical interpretation and small‑discordant behavior: The guarded GVR pipeline produced a small absolute improvement (\textasciitilde{}1.88 percentage points). Under the preregistered exact McNemar decision rule the improvement did not reach \ensuremath{\alpha}=0.05 (p=0.0625). The paired bootstrap CI that excludes zero highlights finite‑sample discreteness: with few discordants (n10=5, n01=0) McNemar's discrete two‑sided p‑value can be conservative for small effects while bootstrap percentile CIs reflect resampled paired‑difference variability. We report both inferential perspectives so readers may weigh formal hypothesis rejection against effect magnitude and CI width.

Operational implications and boundary conditions grounded in artifacts: Baseline first‑pass accuracy was high (\ensuremath{\approx}97.7\%), leaving few failures for the repair stage to correct---hence rare repair invocations (6/300). Stage accounting and provider traces show repair-stage overhead was small in the observed run, but we did not archive per‑episode latency and therefore cannot quantify repair latency or end‑to‑end timing against the preregistered operational bound (median <= 60 s). The low repair rate suggests bounded automated repair adds limited marginal benefit on high‑first‑pass corpora; however, GVR may yield larger absolute improvements when first‑pass accuracy is lower, verifier coverage expands, or when the verifier detects failure modes common in production workloads.

Threats to validity (artifact‑grounded): Internal validity: the shared\_initial\_candidate estimand measures the effect of validator‑triggered repair for a fixed initial sample and does not assess independent multiple draws or constrained first‑pass prompting because independent first‑pass arms were not executed under the preregistration. Construct validity: deterministic\_netops\_validator\_v5 enforces a finite, archived rule battery (syntactic parsing, required\_sequence, constrained\_field checks, reachability) and therefore semantic violations outside that battery are undetected and unmeasured. External validity: tasks derive from public NetConfEval‑style examples and a seeded synthetic generator limited to two vendor dialects; results may not generalize to other vendors, larger topologies, or production heterogeneous workloads. Conclusion validity: the loss of 34 both‑arm provider failures reduced effective N and the observed few discordant pairs limit power for small effect detection.

Design lessons and recommendations grounded in execution artifacts: (1) When baseline pass rates are high, design experiments targeting lower‑baseline corpora or include arms that modify first‑pass sampling (e.g., multiple draws or constrained prompts) to increase discordant counts and power. (2) Archive per‑episode pipeline timing (request\ensuremath{\rightarrow}generation\ensuremath{\rightarrow}validation\ensuremath{\rightarrow}repair\ensuremath{\rightarrow}final validation) to assess operational latency bounds; our run lacked per‑episode latency and thus the preregistered latency estimand could not be evaluated. (3) Expand and publish validator rule coverage and include tests that exercise verifier false negatives; per‑episode validator traces permit such audits and should be part of future preregistrations. (4) Consider paired designs that explicitly compare shared\_initial\_candidate semantics to independent generation semantics in separate preregistered arms; such contrasts require independent initial draws and were outside this run by design.

Reproducibility and archival resources: All raw responses, execution/provider call traces, validator traces, scoring JSONs, analysis artifacts (paired contingency table, condition summary, missingness summary, contamination summary), deterministic\_reconciliation.json, and model\_configuration.json are archived and hashed in the evidence bundle. The archived analysis\_log.jsonl records bootstrap settings (B=10,000, seed=1729). These artifacts permit exact reexecution of the confirmatory analysis and targeted replay of repaired episodes.

\section{Limitations}
\begin{itemize}
\item Shared\_initial\_candidate semantics: both arms used the same single sampled initial candidate per preregistration; the estimand measures validator‑triggered repair benefit for that fixed initial sample and does not evaluate independent first‑pass generations or constrained first‑pass prompting strategies.
\item Missingness and sample reduction: planned N=300 pairs; after infrastructure/provider exclusions 266 pairs were complete and 34 both‑arm episodes were excluded per preregistered rules, reducing effective sample size and precision.
\item Small discordant counts: primary inference used exact McNemar with few discordants (n10=5, n01=0); conclusions are sensitive to inferential choice and limited power.
\item Validator coverage: deterministic\_netops\_validator\_v5 enforces a finite set of syntactic/semantic/reachability rules; semantic violations outside its coverage are possible and unmeasured.
\item Runtime nondeterminism and sampling: archived execution/model\_configuration.json records temperature = null and no deterministic seed; nondeterministic sampling cannot be excluded for recorded model calls.
\item H2 latency not evaluated: per‑episode end‑to‑end latency was not present in archived per‑task artifacts and therefore the preregistered latency bound (median <= 60 s) could not be assessed.
\item External validity: task corpus derives from public NetConfEval‑style examples and a seeded synthetic generator limited to two vendor‑dialect clusters; results do not imply universal benefit across other vendors, larger topologies, or production deployments.
\end{itemize}

\section{Conclusion}
We executed a preregistered paired experiment that evaluated a GVR pipeline (gpt‑5‑mini + deterministic\_netops\_validator\_v5 + <=1 automated repair) on intent\ensuremath{\rightarrow}CLI tasks under shared\_initial\_candidate semantics. Across 266 analysed pairs the guarded pipeline improved validator pass rate by 0.01879699248 but did not meet the preregistered exact‑McNemar \ensuremath{\alpha}=0.05 decision rule (p=0.0625); a paired bootstrap CI (B=10,000, seed=1729) narrowly excludes zero. Repair calls were rare (6/300) and median repair iterations per task = 0. Per‑episode latency was not archived and could not be assessed. All execution and analysis artifacts (raw responses, validator traces, provider call accounting, and reconciliation manifests) are archived in the evidence bundle to support reanalysis and follow‑on work.

References
[1] C. Wang, M. Scazzariello, A. Farshin, S. Ferlin, D. Kostić, and M. Chiesa, "NetConfEval: Can LLMs Facilitate Network Configuration?," in Proc. ACM, 2024, doi: 10.1145/3656296.
[2] S. Y. Long, J. Tan, B. Mao, F. Tang, Y. Li, M. Zhao, and N. Kato, "A Survey on Intelligent Network Operations and Performance Optimization Based on Large Language Models," IEEE Commun. Surveys \& Tutorials, 2025, doi: 10.1109/comst.2025.3526606.
[3] R. Cadenas, "Convergent Correctness in Stochastic Code Generation: A Generate‑Verify‑Repair Architecture for Deterministic Validation of Autonomous AI Coding Agents in Regulated Environments," SSRN, 2026, doi: 10.2139/ssrn.6754899.
[4] M. Brown, A. Fogel, D. Halperin, V. Heorhiadi, R. Mahajan, and T. Millstein, "Lessons from the evolution of the Batfish configuration analysis tool," Proc., 2023, doi: 10.1145/3603269.3604866.
[5] R. Mondal, A. Tang, R. Beckett, T. Millstein, and G. Varghese, "What do LLMs need to Synthesize Correct Router Configurations?," Proc., 2023, doi: 10.1145/3626111.3628194.

\section*{Disclosure Statement}
Disclosure (concise): Declared model: gpt-5-mini (execution/model\_configuration.json archived; declared\_model\_version = "gpt-5-mini"). Orchestration adapter: hosted\_netops\_gvr\_v1. Deterministic validator: deterministic\_netops\_validator\_v5. Preregistration: cnsms2026-001-gvr-batfish-prereg-v1 (manifest SHA‑256: 46a22349\allowbreak{}2f760950\allowbreak{}b3d06f47\allowbreak{}5129ce21\allowbreak{}e52e296c\allowbreak{}4c2a1e09\allowbreak{}df3ccd50\allowbreak{}bec9f891). Initial master prompt immutable reference: provenance/master\_prompt.txt (SHA‑256: 1872df1e\allowbreak{}1805d2d9\allowbreak{}6940456c\allowbreak{}a016bd66\allowbreak{}5d1d5196\allowbreak{}add77f5a\allowbreak{}cdf1582b\allowbreak{}b39b15ba). Human role and autonomy boundary: before the autonomy lock a human Research Director selected the topic family and approved pre‑lock infrastructure development; after the autonomy lock the registered pipeline autonomously performed literature selection, preregistration, experimental execution, analysis, manuscript drafting, AI peer review simulation, and revision. Nondeterminism disclosure: archived execution/model\_configuration.json records temperature = null (unset) and no deterministic sampling seed is archived; therefore nondeterministic sampling of model outputs cannot be excluded for the recorded provider calls. Execution and analysis artifacts, logs, and hash manifests are archived in the evidence bundle referenced in the preregistration.

\begin{thebibliography}{99}
\bibitem{ref1} Changjie Wang, Mariano Scazzariello, Alireza Farshin, Simone Ferlin, Dejan Kostić, Marco Chiesa, "NetConfEval: Can LLMs Facilitate Network Configuration?", 2024, 10.1145/3656296
\bibitem{ref2} S.Y. Long, Jingjing Tan, Bomin Mao, Fengxiao Tang, Yangfan Li, Ming Zhao, Nei Kato, "A Survey on Intelligent Network Operations and Performance Optimization Based on Large Language Models", 2025, 10.1109/comst.2025.3526606
\bibitem{ref3} Rafael Cadenas, "Convergent Correctness in Stochastic Code Generation: A Generate-Verify-Repair Architecture for Deterministic Validation of Autonomous AI Coding Agents in Regulated Environments", 2026, 10.2139/ssrn.6754899
\bibitem{ref4} Matt Brown, Ari Fogel, Daniel Halperin, Victor Heorhiadi, Ratul Mahajan, Todd Millstein, "Lessons from the evolution of the Batfish configuration analysis tool", 2023, 10.1145/3603269.3604866
\bibitem{ref5} Rajdeep Mondal, Alan Tang, Ryan Beckett, Todd Millstein, George Varghese, "What do LLMs need to Synthesize Correct Router Configurations?", 2023, 10.1145/3626111.3628194
\end{thebibliography}

\end{document}
